\subsection{李群中的换位子}

命题 3. 设 G 是有限维李群，\(\mathbf{H}_1\) 和 \(\mathbf{H}_2\) 是 G 的子群。设 \(\mathfrak{h}_1, \mathfrak{h}_2\) 和 \(\mathfrak{h}\) 分别是 \(e\)\(\mathbf{H}_1\mathbf{H}_2\) 和 \((\mathbf{H}_1,\mathbf{H}_2)\) 在 e 处相切的李子代数。则 \([\mathfrak{h}_1,\mathfrak{h}_2] \subset \mathfrak{h}\)。

令 \(a \in \mathfrak{h}_1\)、\(b \in \mathfrak{h}_2\)。存在 K 中 0 的一个开邻域 I 以及从 I 到 G 的解析映射 \(f_1, f_2\)，使得

\[
\begin{array}{c} f _ {1} (0) = f _ {2} (0) = e, \quad f _ {1} (\mathrm{I}) \subset \mathrm{H} _ {1}, \quad f _ {2} (\mathrm{I}) \subset \mathrm{H} _ {2}, \\ (\mathrm{T} _ {0} f _ {1}) 1 = a, \quad (\mathrm{T} _ {0} f _ {2}) 1 = b. \end{array}
\]

记

\[
f (\lambda , \mu) = (f _ {1} (\lambda), f _ {2} (\mu)) \in (H _ {1}, H _ {2}) \quad \text { for   } \lambda , \mu \text {   in   I. }
\]

借助一个将 e 映到 0 的图表，把 G 中 e 的一个开邻域与 \(e\)\(G\)\(K^r\) 的一个开子集识别。于是 \(e\)\(L(G)\) 被识别为 \(K^r\)。由 §5，第 2 号，命题 1，\(f(\lambda, \mu)\) 关于原点的整级数展开为

\[
f (\lambda , \mu) = \lambda \mu [ a, b ] + \sum_ {i \geqslant 1, j \geqslant 1, i + j \geqslant 2} \lambda^ {i} \mu^ {j} a _ {i j}
\]

其中 \(a_{ij} \in \mathbf{K}^r\)（因为  的展开式中 \(\lambda^i\) 和 \(\mu^j\) 的项为零，这是由于 \(f(\lambda, \mu)\)\(f(\lambda, 0) = f(0, \mu) = 0\)）。固定 I 中的 \(\mu\)。令 \(\lambda\) 趋于 0，可见

\[
\mu [ a, b ] + \sum_ {j \geqslant 2} \mu^ {j} a _ {1 j} \in \mathfrak {h}.
\]

由于这对所有 \(\mu\in\mathbf{I}\) 都成立，便有 \([a,b]\in\mathfrak{h}\)。

注记。即使 H＿1 和 H＿2 是 G 的连通李子群，由 \(H_{1}\)\(H_{2}\)\([h_{1}, h_{2}]\) 生成的 L(G) 的李子代数一般也不同于 h。

命题 4. 设 \(\mathbf{G}\) 是有限维实或复李群，\(\mathbf{A}, \mathbf{B}, \mathbf{C}\) 是 \(\mathbf{G}\) 的积分子群，满足 \([\mathrm{L}(\mathrm{A}), \mathrm{L}(\mathrm{C})] \subset \mathrm{L}(\mathrm{C})\) 且

\[
[ \mathrm{L} (\mathrm{B}), \mathrm{L} (\mathrm{C}) ] \subset \mathrm{L} (\mathrm{C}).
\]

若 \([\mathrm{L}(\mathrm{A}),\mathrm{L}(\mathrm{B})]\subset \mathrm{L}(\mathrm{C})\)，则 \((\mathrm{A},\mathrm{B})\subset \mathrm{C}\)。若 \([\mathrm{L}(\mathrm{A}),\mathrm{L}(\mathrm{B})] = \mathrm{L}(\mathrm{C})\)，则 \((\mathrm{A},\mathrm{B}) = \mathrm{C}\)。

假设 \([\mathrm{L}(\mathrm{A}),\mathrm{L}(\mathrm{B})]\subset \mathrm{L}(\mathrm{C})\)。\(\mathrm{L}(\mathrm{A}) + \mathrm{L}(\mathrm{B}) + \mathrm{L}(\mathrm{C})\) 是 \(\mathrm{L}(\mathrm{G})\) 的李子代数。取 G 中李代数为 \(\mathbf{G}\)\(\mathrm{L}(\mathrm{A}) + \mathrm{L}(\mathrm{B}) + \mathrm{L}(\mathrm{C})\) 的积分子群后，问题归结为

\[
\mathrm{L} (\mathrm{A}) + \mathrm{L} (\mathrm{B}) + \mathrm{L} (\mathrm{C}) = \mathrm{L} (\mathrm{G})
\]

且 G 连通。因此 L(C) 是 L(G) 的理想。先假设 G 单连通。则 C 是 G 的正规李子群（§ 6，第 6 号，命题 14）。令 \(\S\)\(\phi\) 为从 G 到 G/C 的典范态射。于是

\[
[ \mathrm{L} (\phi) (\mathrm{L} (\mathrm{A})), \mathrm{L} (\phi) (\mathrm{L} (\mathrm{B})) ] = \{0 \}
\]

从而由 Hausdorff 公式，\(\phi (\mathbf{A})\) 与 \(\phi (\mathbf{B})\) 交换；因此 \((\mathrm{A},\mathrm{B})\subset \mathrm{C}\)。一般地，令 \(\mathbf{G}'\) 为 G 的万有覆盖，令 \(\mathbf{A}',\mathbf{B}',\mathbf{C}'\) 为 \(\mathbf{G}'\) 中满足 \(\mathrm{L}(\mathrm{A}') = \mathrm{L}(\mathrm{A}),\mathrm{L}(\mathrm{B}') = \mathrm{L}(\mathrm{B}),\)、、\(\mathrm{L}(\mathrm{C}') = \mathrm{L}(\mathrm{C})\) 的积分子群。于是 \((\mathrm{A}',\mathrm{B}')\subset \mathrm{C}'\)，而 A、B、C 是 A'、B'、C' 在 G 中的典范像，从而 \((\mathrm{A},\mathrm{B})\subset \mathrm{C}\)。另一方面，(A, B) 是 G 的某个积分子群的底层集合（§ 6，第 2 号，命题 4 的推论），其李代数包含 [L(A), L(B)]（命题 3）。若

\[
[ \mathrm{L} (\mathrm{A}), \mathrm{L} (\mathrm{B}) ] = \mathrm{L} (\mathrm{C}),
\]

则 (A, B) ⊃ C，从而 (A, B) = C。

推论。设 G 是带有李代数 g 的有限维连通实或复李群。子群 D \(^{i}\) G（分别地，C \(^{i}\) G）是积分子群，其李代数为 D \(^{i}\) g（分别地，C \(^{i}\) g）。若 G 单连通，则它们是李子群。

第一个断言由命题 4 对 i 作归纳得到。第二个断言由第一个断言和 § 6，第 6 号，命题 14 得到。

命题 5. 设 G 是有限维实或复李群，A 是 G 的积分子群。则 \(\overline{\mathrm{D}\mathrm{A}} = \mathrm{D}\overline{\mathrm{A}}\)。特别地，A 是 \(\overline{\mathbf{A}}\) 的正规子群，且 \(\overline{\mathbf{A}} / \mathbf{A}\) 是交换群。

令 \(\mathfrak{a} = \mathrm{L}(\mathbf{A})\)。令 \(G_{1}\) 为满足下列条件的 \(g \in G\) 的集合：

\[
(\operatorname{Ad} g) x \equiv x (\mathrm{mod} \mathscr {D} \mathfrak {a}) \quad \text { for   all } x \in \mathfrak {a}.
\]

于是 \(G_1\) 是 G 的闭子群。若 \(G\)\(y \in \mathfrak{a}\)，则由 § 6，第 4 号，命题 10 的推论 3 (ii)，\(\exp y \in G_1\)。因此 \(G_1\) 包含 A，进而包含 \(A\)\(\bar{A}\)。所以对于 \(g \in \bar{A}\)，\(L(\text{Int } g)\) 保持 a 不变，因而 \(a\)\(\text{Int } g\) 保持 A 不变；更确切地说，\(A\)\(L(\text{Int } g)\) 在 \(a / \mathcal{D}a\) 上定义恒等自同构，因而 \(\text{Int } g\) 在 \(A / DA\) 上定义恒等自同构。这证明了 \((\bar{A}, A) \subset DA\)。赋予 G 实李群结构后，\(G\)\(\bar{A}\) 是李子群（§ 8，第 2 号，定理 2）；令 b 为其李代数。令 \(\S 8\)\(b\)\(G_2\) 为满足下列条件的 \(g \in G\) 的集合：

\[
(\operatorname{Ad} g) x \equiv x (\mathrm{mod} \mathscr {D} \mathfrak {a}) \quad \text { for   all } x \in \mathfrak {b}.
\]

由上可知，\(\mathbf{G}_2\supset \mathbf{A}\)，因而 \(\mathbf{G}_2\supset \overline{\mathbf{A}}\)。所以对于 \(g\in \overline{\mathbf{A}}\)，Int \(g\) 保持 \(\mathrm{D}\overline{\mathbf{A}}\) 不变，并在 \(\overline{\mathbf{A}} / \mathrm{D}\overline{\mathbf{A}}\) 上定义恒等自同构。因此 \(\mathrm{D}\overline{\mathbf{A}}\supset \overline{\mathrm{D}\mathbf{A}}\)。

命题 6. 假设 K 是超度量的。设 G 是有限维李群，A、B、C 是 G 的李子群，满足 \([L(A), L(C)] \subset L(C)\)、\([L(B), L(C)] \subset L(C)\)。若 \([L(A), L(B)] \subset L(C)\)，则存在 A、B 的开子群 \(A'\)、\(B'\)，使得 \((A', B') \subset C\)。若 \([L(A), L(B)] \subset L(C)\)，则存在 A、B、C 的开子群 \(A'\)、\(B'\)、\(C'\)，使得 \((A', B') = C'\)。

假设 \([\mathrm{L}(\mathrm{A}),\mathrm{L}(\mathrm{B})]\subset\mathrm{L}(\mathrm{C})\)。如命题 4 的证明，问题归结为 \(\mathrm{L}(\mathrm{C})\) 是 \(\mathrm{L}(\mathrm{G})\) 的理想的情形。然后，通过将 G 替换为一个开子群，问题归结为 C 是 G 的正规子群的情形（§ 7，第 1 号，命题 2）。令 \(\S\)\(\phi\) 为从 G 到 G/C 的典范态射。于是

\[
[ \mathrm{L} (\phi) (\mathrm{L} (\mathrm{A})), \mathrm{L} (\phi) (\mathrm{L} (\mathrm{B})) ] = \{0 \}.
\]

由 Hausdorff 公式，存在 A、B 的开子群 \(\mathbf{A}'\)、\(\mathbf{B}'\)，使得 \(\mathbf{A}\)\(\mathbf{B}\)\(\phi(\mathbf{A}')\) 与 \(\phi(\mathbf{B}')\) 交换，从而 \((\mathbf{A}', \mathbf{B}') \subset \mathbf{C}\)。再假设

\[
[ \mathrm{L} (\mathrm{A}), \mathrm{L} (\mathrm{B}) ] = \mathrm{L} (\mathrm{C}).
\]

由命题 3，在 e 处相切于 \((\mathbf{A}',\mathbf{B}')\) 的李子代数包含 \(e\)\(\mathrm{L}(\mathrm{C})\)。因此 \((\mathbf{A}',\mathbf{B}')\) 包含一个李代数为 \(\mathbf{G}\)\(\mathrm{L}(\mathrm{C})\) 的 G 的李子群芽。所以 \((\mathbf{A}',\mathbf{B}')\) 是 \(\mathbf{C}\) 的开子群。

推论。假设 K 是超度量的。设 G 是李代数为 g 的有限维李群。存在 G 的一个开子群 \(G_{0}\)，使得对所有 i，\(D^{i}G_{0}\)（分别地，\(C^{i}G_{0}\)）是 G 的李子群，其李代数为 \(D^{i}g\)（分别地，\(C^{i}g\)）。

\begin{enumerate}
\def\labelenumi{(\alph{enumi})}
\item
  由命题 3 归纳应用可知，对于 G 的每个开子群 \(G_{1}\) 和所有 i，\(D^{i}G_{1}\) 都包含一个李代数为 \(D^{i}g\) 的 G 的李子群芽。
\item
  设 \(G'\)\(G\)\(i \leqslant n\) 是 G 的开子群，使得当 \(D^i G'\) 时，\(G\)\(D^i\mathfrak{g}\) 是 G 的李子群，其李代数为 \(H_1, H_2\)\(D^n G'\)。由命题 6，存在 \((H_1, H_2)\)\(D^{n+1}\mathfrak{g}\) 的开子群 ，使得  是李子群，其李代数为 \(G''\)。令 \(G'\) 为 G' 的开子群，取得足够小，使得 \(D^n G'' \subset H_1 \cap H_2\)。则 \(D^{n+1}G'' \subset (H_1, H_2)\)。这些关系
\end{enumerate}

\[
\mathrm{D} ^ {0} \mathrm{G} ^ {\prime \prime} \subset \mathrm{D} ^ {0} \mathrm{G} ^ {\prime}, \mathrm{D} ^ {1} \mathrm{G} ^ {\prime \prime} \subset \mathrm{D} ^ {1} \mathrm{G} ^ {\prime}, \dots , \mathrm{D} ^ {n} \mathrm{G} ^ {\prime \prime} \subset \mathrm{D} ^ {n} \mathrm{G} ^ {\prime}, \mathrm{D} ^ {n + 1} \mathrm{G} ^ {\prime \prime} \subset (\mathrm{H} _ {1}, \mathrm{H} _ {2})
\]

由 (a) 可知，对于 \(\mathbf{D}^i\mathbf{G}''\)，\(i\leqslant n + 1\) 是 \(\mathbf{G}\) 的李子群，其李代数为 \(\mathcal{D}^i\mathfrak{g}\)。

\begin{enumerate}
\def\labelenumi{(\alph{enumi})}
\setcounter{enumi}{2}
\item
  存在整数 p，使得 \(D^{p}g = D^{p+1}g = \cdots\)。由上可知，存在 G 的开子群 \(G_{0}\)，使得当 \(D^{i}G_{0}\) 时，\(i \leqslant p\) 是 G 的李子群，其李代数为 \(D^{i}\mathfrak{g}\)。但是由 (a)，当 \(i > p\)\(D^{p}G_{0} \supset D^{i}G_{0}\) 时同一断言仍成立，因为 \(i > p\)。
\item
  对 \(C^{i}\) 的论证类似。
\end{enumerate}

